\documentclass[book.tex]{subfiles}
\begin{document}

\chapter{Neural Equivariant Potentials}
\label{chap:neural_equivariant_potentials}
\noindent\rule{11cm}{0.4pt} \\
\textbf{Prerequisites:}  \autoref{chap:molecular_hamiltonian},   \\
\textbf{Difficulty Level:} ** \\
\noindent\rule{11cm}{0.4pt}
\vspace{5mm}

%"""
%Representations of atomistic systems for machine learning must transform predictably under the geometric transformations of 3D space, in particular rotation, translation, mirrors, and permutation of atoms of the same species. These constraints are typically satisfied by means of atomistic representations that depend on scalar distances and angles, leaving the representation invariant under the above transformations. Invariance, however, limits the expressivity and can lead to an incompleteness of  representations. In order to overcome this shortcoming, we recently introduced Neural Equivariant Interatomic Potentials (NequIP) [1], a Message Passing Neural Network approach for learning interatomic potentials that uses an E(3)-equivariant representation of atomic environments. While most current interatomic potentials based on Message Passing Neural Network use invariant convolutions over scalar features, NequIP instead employs equivariant convolutions over geometric tensors (scalar, vectors, …), providing a more information-rich message passing scheme. In my talk, I will first motivate the choice of an equivariant representation for atomistic systems and demonstrate how it allows for the design of interatomic potentials at previously unattainable accuracy. I will discuss applications on a diverse set of molecules and materials, including small organic molecules, water in different phases, a catalytic surface reaction, glass formation of a lithium phosphate, Li diffusion in a superionic conductor, as well as protein folding. I will then show that NequIP can predict structural and kinetic properties from molecular dynamics simulations in excellent agreement with ab-initio simulations. The talk will then discuss the observation of a remarkable sample efficiency in equivariant interatomic potentials which outperform existing neural network potentials with up to 1000x fewer training data and rival or even surpass the sample efficiency of kernel methods. Finally, I will discuss potential reasons for the high sample efficiency of equivariant interatomic potentials.
%"""
%References .

A broad goal of molecular modeling is to derive a machine learned potential that is as accurate as an ab initio potential but as fast as a classical potential. Machine learned potentials have the powerful advantage of having linear cost in the number of atoms $n$ and with the potential for ab initio accuracy at near classical speeds due to the power of the universal approximation theorem.

In previous chapters, we have studied various techniques for representing atomic systems for machine learned potentials. One emerging idea is that representations of atomic system for machine learned potentials should transform correctly under transformations of three dimensional space such as rotation, translation, reflection and identical atom permutation. In this chapter, we introduce an equivariant architecture for predicting properties of molecular and atomic systems.

\section{Invariances and Equivariances in Molecular Architectures}

Recall that $E(3)$ is the group of isometries in Euclidean space $\mathbb{R}^3$ that leave distances constant. Note that any function $f$ that depends only on $\|r_i - r_j\|$ inputs is $E(3)$ invariant. Thus the following function is invariant to $E(3)$

\begin{align}
    f(r_1-r_2, r_1-r_3,\dotsc)
\end{align}

This fact has been used in classical potentials and not just for machine learned potentials. The limitation of invariance in this context is that it only applies to scalar quantities. To handle vector and tensor properties, we can extend invariance to equivariance. Recall that formally, a function $f: X \to Y$ is equivariant to group $G$ that acts on both $X$ and $Y$ if for all $g \in G$, $x \in X$ we have

\begin{align}
f(D_X(g)x) = D_Y(g)f(x)  
\end{align}

where $f: X \to Y$ is a vector space map. In our case $G = E(3)$. The Behler-Parinello representations we saw earlier are almost equivariant, but the interatomic cutoffs imposed in that architecture break the equivariance. %(\textcolor{red}{TODO: Double check this statement and prove})

\section{Constructing Invariant Representations}
In this section, we discuss how to build a general family of $E(3)$ invariant operators. Invariance is constructed by composing a number of properties:

\begin{itemize}
    \item Translation: Operate on interatomic distances $r_i - r_j$ to enforce translation invariance.
    \item Permutation: Summation is permutation invariant, so the computed atomic feature vector sums over neighboring atoms vectors.
    \item Rotational: Constrain filters to be $SO(3)$ invariant.
    \item Parity: Add a output contribution rule that respects parity.
\end{itemize}

\section{Features and Convolutions}

At a high level, NequIP follows the Behler-Parinello atomic energy contribution paradigm
\begin{align}
    E &= \sum_{i} E_i
\end{align}
where $E_i$ are the atomic energy contributions for atom $i$. Every atom is associated with scalar, vector and tensor features. Feature vectors here are geometric objects that are direct sums of irreducible representations of $O(3)$ %(\textcolor{red}{TODO: Isn't an irreducible representation a function? How is that transformed into a concrete geometric object}). 
The feature vectors are given by
\begin{align}
    V^{l,p}_{acm}
\end{align}
where $l = 0, 1, \dotsc$ is the rotation order, and $p \in \{+1,-1\}$ is the parity. $l, p$ together label the relevant irreducble representation of $O(3)$. $a$ is the atom, $c$ is the channel in the feature vector, and $m$ is the representation index in $[-l, l]$. Convolutions are performed on these geometric objects and are defined by

\begin{align}
F(\vec{r}_{ij}) &= R(r_{ij})Y^{(l)}_m(\hat{r}_{ij})
\end{align}
where $Y^l_m$ is a spherical harmonic function (equivariant under $SO(3)$ %\textcolor{red}{TODO: Add a proof of this earlier}). 
Here 
\begin{align}
\vec{r}_{ij} &= \vec{r}_i - \vec{r}_j \\
\hat{r}_{ij} &= \frac{\vec{r}_{ij}}{\|\vec{r}_{ij}\|} \\
r_{ij} &= \|\vec{r}_{ij}\|
\end{align}
$R$ is a rotationally invariant function given by
\begin{align}
    R(r_{ij}) &= W_n\sigma(\dotsc \sigma(W_2\sigma(W_1 B(r_{ij}))))
\end{align}

Here $W_i$ are weight matrices, $\sigma$ is the nonlinear activation, and $B$ is defined as

\begin{align}
    B(r_{ij}) &= \frac{2}{r_c} \frac{\sin(\frac{n\pi}{r_c})}{r_{ij}}f_{env}(r_{ij}, r_c)
\end{align}
where $r_c$ is a radial cutoff distance 
%(\textcolor{red}{NOTE: This means our explanation earlier about BP cutoffs is wrong. Fix}) 
and $f_{env}$ is a polynomial 
%(\textcolor{red}{TODO: Find a concise description of the polynomial here}). 
Note that the use of $r_{ij}$ to compute all features provides translational invariance. Parity is enforced by placing a constraint
\begin{align}
    p_o &= p_i p_f
\end{align}

The full convolution is given by

\begin{align}
o: l_i \otimes l_f &\to l_o \\
L^{(l_o, p_o, l_f, l_i, p_f, p_i)}_{acm_o}(\vec{r}_a, V^{(l_i, p_i)}_{acm}) &= \sum_{m_f, m_i} C^{(l_o, m_o)}_{(l_i, m_i), (l_f, m_f)} \sum_{b \in S} R^{(l_f, l_i, p_f, p_i)}_c(r_{ab})Y^{(l_f)}_{m_f}(\hat{r}_{ab}) V^{(l_i, p_i)}_{bcm_i}
\end{align}
%\textcolor{red}{TODO: Explain this convolution in more detail.}

Features and filters are combined through equivariant tensor product.
\begin{align}
    D_X \otimes D_Y
\end{align}
%(\textcolor{red}{TODO: Unify notation between here and the Behler-Parinello Chapter}

Message passing neural networks have emerged as a powerful type of network for neural potentials. Let $M_t$ be a message function. $h^t_v$ is a vector at node $v$.

\begin{align}
    m^{t+1}_v &= \sum_{w \in N(v)} M_t(h^t_v, h^t_{w}, e_{vw}) \\
    h^{t+1}_v &= U(h^t_v, m^{t+1}_v)
\end{align}

\section{Equivariant Architecture}
\begin{figure}
    \centering
    % \includegraphics{figures/Differentiable Physics/Neural Potentials/NequIP/nequip_architecture.png}
    \includegraphics{figures/Differentiable Physics/Neural Potentials/NequIP/NequIP architecture.png}
    \caption{A diagram of the NequIP architecture.}
    \label{fig:nequip_architecture}
\end{figure}
NequIP\cite{batzner2021se} is an $E(3)$ equivariant message passing network. The difference between NequIP and SchNet is that atoms don't only carry feature vectors of scalars but also vector and tensor features. Exploiting equivariance on these higher geometric objects allows for richer learned representations.

In addition, convolutions are not invariant to $E(3)$ but are equivariant using a functional form from tensor field networks.

%Equivariance is by $x$



%E3NN introduces these ideas are pseudoscalars, pseudovectors, pseudotensors E3NN features are indexed by irreducible representations of $O(3)$.

%$nxlp$. $p$ is parity. $n$ is multiplicity. $l$ is rotation order.$m$ is rotation index. Use spherical harmonics $Y^{(l)}_m$.

%Notes that SchNet, DIMENet are far more efficient than DeepMD. Much mmore sample efficient. Note that NequIP is massively overparameterized but does quite well.

%NequIP authors claimm neural potential no longer need giant training datsets.

%They are able to use NequIP on amorphous materials structures as well. Able to apply on a dataset of protein fragments (comparison against PhysNet). NequIP also generalizes to high temperature MD. Also complex dihedrals.

\section{Sample Efficiency}

A powerful advantage of leveraging equivariances is increased sample efficiency. Empirically, equivariant architectures like NequIP are able to learn with orders of magnitude less data than simpler architectures in some cases.

One idea for the sample efficiency is that the fact that internal features are vectors means that richer internal representations can be exploited.

The resulting architectures also appear to be more generalizable, with demonstrated applicability to problems in amorphous materials characterization, protein characterizations, dihedral calcluations and more. %(\textcolor{red}{TODO: CITE})
\end{document}